\documentclass[usenames,dvipsnames]{beamer}

% French characters
\usepackage[utf8x]{inputenc}
\usepackage[T1]{fontenc}
% French display
\usepackage[english]{babel}
% Math package
\usepackage{amsmath}
\usepackage{amsfonts}
\usepackage{amssymb}
% includegraphics
\usepackage{graphicx}
% units
%\usepackage{hepunits}
% links
\usepackage{hyperref}
% \usepackage[bookmarksopen,bookmarksopenlevel=1]{hyperref}
\usepackage{colortbl}
\usepackage{tikz}
\usepackage{rotating}
% Theme & Color theme
\usepackage{default}
\usetheme{Frankfurt}
\usecolortheme{perso}
% My green
\definecolor{MyColor}{RGB}{43,153,43}
\setbeamercolor{structure}{fg=MyColor}
% \setbeamercolor{titlelike}{fg=MyColor}
% \setbeamercolor{palette primary}{fg=MyColor}
% \setbeamercolor{palette secondary}{fg=MyColor}
% \setbeamercolor{palette tertiary}{bg=MyColor}
% \setbeamercolor{palette quaternary}{bg=MyColor}
% Font
\usepackage{times}
\usefonttheme[onlymath]{serif}
% Suppress the navigation bar
% \beamertemplatenavigationsymbolsempty
\setbeamertemplate{navigation symbols}{}

% Outerthemes
% \useoutertheme[footline=empty, subsection=false]{miniframes}
% \useoutertheme[subsection=false]{miniframes}
% \useoutertheme[subsection=false]{smoothbars}
% gets rid of bottom navigation bars
\setbeamertemplate{footline}[page number]{}
% 
% \setbeamercolor{frametitle}{parent=palette primary}
% \setbeamercolor{subsection in head/foot}{parent=palette secondary}
% \setbeamercolor{section in head}{parent=palette quaternary}
% \setbeamercolor{section in sidebar}{fg=black}
% \setbeamercolor{section in sidebar shaded}{fg= grey}
% Appendix page numbering
\usepackage{appendixnumberbeamer}
% feynmp for Feynman diagrams
%\usepackage{feynmp}
%\DeclareGraphicsRule{*}{mps}{*}{}
\makeatletter
\def\endfmffile{%
  \fmfcmd{\p@rcent\space the end.^^J%
          end.^^J%
          endinput;}%
  \if@fmfio
    \immediate\closeout\@outfmf
  \fi
  \ifnum\pdfshellescape=\@ne
    \immediate\write18{mpost \thefmffile}%
  \fi}
\makeatother
\usepackage[abs]{overpic}
\def\met{\ensuremath{\not\!\!{E_{T}}}}

\newcommand{\ttbar}{$t\bar{t}$}
\newcommand{\bbbar}{$b\bar{b}$}
\newcommand{\GeVcc}{GeV/$\text{c}^2$}
\newcommand{\pt}{$p_{T}$}
\newcommand{\ptg}[1]{$p_{T}>#1$~GeV/c}
\newcommand{\etal}[1]{$|\eta|<#1$}
\newcommand{\etag}[1]{$|\eta|>#1$}
\newcommand{\HT}{$H_{T}$}
\newcommand{\MET}{$\slashed{E}_{T}$}
\newcommand{\METg}[1]{$\slashed{E}_{T}>#1$~GeV/c}
\newcommand{\HTg}[1]{${H_{T}>#1}$~GeV/c}
\newcommand{\jn}[1]{$j_{#1}$}
\newcommand{\ex}[1]{\times 10^{#1}}
\newcommand{\Z}{$Z^{0}$}
\newcommand{\tq}{$t$}
\newcommand{\W}{$W^{+/-}$}
\newcommand{\Hb}{$H^{0}$}
\newcommand{\Tp}{$T'$}
\newcommand{\Tj}{$T'j$}
\newcommand{\Tjj}{$T'jj$}

%% Title and such
\title{\texorpdfstring{Search for a vector-like quark \Tp~decaying \\into top+Higgs in single production mode \\in full hadronic final state using \\CMS data collected at 8 TeV\\}}
% \subtitle{
% }
\author[J.~Ruiz-\'{A}lvarez]{Jos\'{e} D. Ruiz-\'{A}lvarez}
\institute[IPN Lyon]{
% Institut de Physique Nucléaire de Lyon / Université Claude Bernard Lyon 1\\
\includegraphics[height=.6cm]{../Prelude/CNRS.png}\hspace{.5cm}
\includegraphics[height=.6cm]{udl.png}\hspace{.5cm}
\includegraphics[height=.6cm]{../Prelude/UCBL.png}\hspace{.5cm}
\includegraphics[height=.6cm]{EDPhast.jpg}
\\\vspace{.5cm}
\includegraphics[height=.7cm]{../Prelude/IPNL.png}\hspace{.5cm}
\includegraphics[height=.7cm]{../Prelude/CMS.png}
}
\date{21 octobre 2015}

% \mode<presentation> {
%   \usefonttheme{serif}
%   \useoutertheme[footline=authorinstitutetitle,subsection=false]{miniframes}
%   \usecolortheme{crane}
%   \useinnertheme[<shadow>]{rounded}
% }

% \setbeamercolor{boum}{bg=\color{white!40}, fg=black}
% 
% \addtobeamertemplate{block begin}{\pgfsetfillopacity{0.65}}{\pgfsetfillopacity{1}}
% 
% \newenvironment<>{varblock}[2][\textwidth]{%
%   \setlength{\textwidth}{#1}
%   \begin{actionenv}#3%
%     \def\insertblocktitle{#2}%
%     \par%
%     \usebeamertemplate{block begin}}
%   {\par%
%     \usebeamertemplate{block end}%
%   \end{actionenv}}
% 

\usepackage{fancybox}
\usepackage{framed,xcolor}
\colorlet{shadecolor}{white!60}
\usepackage{rotating}
% \usepackage[table]{xcolor} 

% \includeonly{11_StandardModel}
% \includeonly{11_StandardModel,12_LHC_CMS}
% \includeonly{12_LHC_CMS}
% \includeonly{11_StandardModel,12_LHC_CMS,13_Zgamma}
% \includeonly{13_Zgamma}
% \includeonly{14_Hgg,15_BSM}
% \includeonly{14_Hgg,BB_Hgg}
% \includeonly{13_Zgamma,14_Hgg,BB_Hgg}
% \includeonly{15_BSM,BB_BSM}
% \includeonly{BB_Zgamma}


\begin{document}

\begin{frame}[plain]
  \tikz[remember picture,overlay]
    \node (b) at
      (current page.center) [inner sep=0pt,opacity=.4]
%       (0.5\paperwidth,-0.5\paperheight) [inner sep=0pt,opacity=.4]
      {\begin{overpic}[height=.9\paperheight]{a-1_16_1_RhoPhi.png}\end{overpic}};
  \titlepage
\end{frame}

\setcounter{tocdepth}{1}
\begin{frame}{Plan}
\tableofcontents
\end{frame}


%\include{11_StandardModel}

%\include{12_LHC_CMS}

%\include{13_Zgamma}

%\include{14_Hgg}

%\include{15_BSM}


%%%%%%%%%%%%%%%%%%%%%%%%%%%%%%%%%%%%%%%%%%%%%%%%%%%%%%%%
%%%%%%%%%%%%%%%%%%%%%%%%%%%%%%%%%%%%%%%%%%%%%%%%%%%%%%%%
\section[Conclusion]{Conclusions, perspectives}
\subsection{Conclusions}
\begin{frame}{Conclusions}
\vspace{-.5cm}
\begin{block}{}
\begin{itemize}\scriptsize
  \item Étude des désintégrations radiatives $Z^0\to\mu\mu\gamma$
  \begin{itemize}\scriptsize
    \item Validation de la simulation des photons
    \begin{itemize}\tiny
      \item Outil mis à disposition de la collaboration
    \end{itemize}
    \item Mesure de l'échelle d'énergie dans les données 2010
    \begin{itemize}\tiny
      \item Document public : CMS-DP-2011-008 (Note interne : CMS-AN-2011-088)
      \item Effet systématique pour la mesure de section efficace $V\gamma$ - PLB 701 (2011) 535-555
    \end{itemize}
    \item Mesure différentielle de l'échelle d'énergie dans les données 2011
    \begin{itemize}\tiny
      \item Document public : https://twiki.cern.ch/twiki/bin/view/CMSPublic/EcalDPGResults (Note interne : CMS-AN-2012-110)
    \end{itemize}
  \end{itemize}
  \item Recherche du boson de Higgs du Modèle Standard dans le canal $H\to\gamma\gamma$
  \begin{itemize}\scriptsize
    \item Participation à l'identification des photons (validation des photons, efficacité de rejet des électrons)
    \item Vérification et étude de la procédure de séparation en classes
    \begin{itemize}\tiny
      \item Publication : Observation of a new boson at a mass of 125 GeV with the CMS experiment at the LHC - PLB 716 (2012) 30-61
      \item Document public : CMS-PAS-HIG-015 (Note interne : CMS-AN-2012-160)
    \end{itemize}
  \end{itemize}
  \item Production associée d'un boson de Higgs dans le cadre d'un modèle de quarks vecteurs
  \begin{itemize}\scriptsize
    \item Géneration avec ALPGEN des bruits de fond pour le canal $h\to\gamma\gamma$
    \item Étude du canal $h\to\gamma\gamma$
    \begin{itemize}\tiny
      \item Proceedings : Les Houches 2011: Physics@TeV Colliders NP WG Report% - arXiv:1203.1488v2
      \item Publication : Physical Review D 85, 115022 (2012) - arXiv:1204.0455v1
    \end{itemize}
  \end{itemize}
\end{itemize}
\end{block}
\end{frame}

\subsection{Perspectives}
\begin{frame}{Perspectives}
\vspace{-.3cm}
\begin{block}{}
\begin{itemize}\scriptsize
  \item Étude des désintégrations radiatives $Z^0\to\mu\mu\gamma$
  \begin{itemize}\scriptsize
    \item Échelle d'énergie pour 
    \item Finaliser l'étude des incertitudes systématiques
  \end{itemize}
  \item Recherche du boson de Higgs du Modèle Standard dans le canal $H\to\gamma\gamma$
  \begin{itemize}\scriptsize
    \item Réoptimisation de la classification utilisée
    \item Observation d'un excès dans le canal diphoton, excès observé dans d'autres canaux (et dans ATLAS) : découverte d'un nouveau boson, compatible avec un boson de Higgs du Modèle Standard
    \item Propriétés de l'excès (spin, couplages, ...) ? Découverte dans le canal $H\to\gamma\gamma$ seul ?
  \end{itemize}
  \item Production associée d'un boson de Higgs dans le canal $h\to\gamma\gamma$ dans le cadre d'un modèle de quarks vecteurs
  \begin{itemize}\scriptsize
    \item Étude prospective en simulation complète
    \item Étude sur les données récoltées à 
  \end{itemize}
\end{itemize}
\end{block}
\vspace{-.1cm}

\begin{block}{}
\begin{itemize}\scriptsize
\item L'analyse porte sur les données 2011 et 2012 et a été rendue publique le 4 juillet dernier. La quantité de données a doublé depuis... et le LHC continue à accumuler des données !
\item Longue période d'arrêt du LHC pour 2013-2014... avant reprise des collisions à 
% \item Caractérisation de la nouvelle particule
\end{itemize}
\end{block}
\end{frame}

\begin{frame}
\vspace{-.5cm}
\begin{center}
Observation of a new boson at a mass of 125 GeV\\with the CMS experiment at the LHC
\\\tiny
PLB 716 (2012) 30-61 - CMS-PAS-HIG-12-020 - arXiv:1207.7235
\end{center}
%\begin{center}
%\includegraphics[width=.4\textwidth]{plots/Higgs/HIG-12-028/fig16.pdf}
%\includegraphics[width=.4\textwidth]{plots/Higgs/HIG-12-028/fig18.pdf}
%\\
%\includegraphics[width=.4\textwidth]{plots/Higgs/HIG-12-028/fig17.pdf}
%\includegraphics[width=.4\textwidth]{plots/Higgs/HIG-12-028/fig19.pdf}
%\end{center}
\vspace{-.8cm}
\begin{center}
Merci de votre attention !
\end{center}
\end{frame}

%%%%%%%%%%%%%%%%%%%%%%%%%%%%%%%%%%%%%%%%%%%%%%%%%%%%%%%%
%%%%%%% BACKUP
%%%%%%%%%%%%%%%%%%%%%%%%%%%%%%%%%%%%%%%%%%%%%%%%%%%%%%%%
\appendix
\section{BACKUP}
\begin{frame}
\begin{center}
\LARGE
BACKUP
\end{center}
\end{frame}

\begin{frame}
\tableofcontents
\end{frame}



\end{document}